% ECONOMICS_PROBLEM: {"id": "E32-502", "title": "Origins of business cycles", "jel_code": "E32", "source_id": "OP-0231"}
% !TeX program = xelatex
\documentclass[11pt,a4paper]{article}
\usepackage[margin=23mm]{geometry}
\usepackage{fontspec}
\setmainfont{Latin Modern Roman}
\usepackage{amsmath,amssymb,mathtools}
\usepackage{xcolor,enumitem,booktabs,longtable,array,xurl,hyperref}
\definecolor{muted}{HTML}{53636C}
\hypersetup{colorlinks=true,urlcolor=blue,linkcolor=blue}
\setlength{\parindent}{0pt}
\setlength{\parskip}{5pt}
\newcommand{\R}{\mathbb R}
\newcommand{\E}{\mathbb E}
\newcommand{\N}{\mathbb N}
\newcommand{\ind}{\mathbf 1}
\newcommand{\Var}{\operatorname{Var}}
\newcommand{\Cov}{\operatorname{Cov}}
\newcommand{\supp}{\operatorname{supp}}
\DeclareMathOperator*{\argmin}{arg\,min}
\DeclareMathOperator*{\argmax}{arg\,max}
\newcommand{\entrymeta}[3]{{\sffamily\small\color{muted}\textbf{\texttt{#1}}\quad|\quad #2\par #3\par}}
\newcommand{\Problemref}[1]{\texttt{#1}}
\newcommand{\JELref}[2]{\texttt{#2} (JEL \texttt{#1})}
\begin{document}
\section*{Origins of business cycles}
\textbf{Problem ID:} \texttt{E32-502}\quad \textbf{JEL:} \texttt{E32}\par
% BEGIN ECONOMICS STATEMENT
\label{problem:OP-0231}
{\sffamily\small\color{muted}Original subject group: Business cycles and macro dynamics\par}
\label{definitions:problem:30007512}
\textbf{Audit classification:} Published research agenda. The source leaves belief-versus-short-run-news interpretation unresolved; the four-way attribution is editorial. Verification concerns the cited version, not first priority or proof that this exact editorial specification remains unsolved on October 8, 2026.
\subsection*{Definitions and precise research target}
\paragraph{Editorial mathematical specification.} Fix a country, a quarterly sample, and a finite vector \(z_t\) containing log output, consumption, investment, hours, prices, and observed financial conditions. Assume a stable reduced-form vector autoregression with innovation covariance \(\Sigma\), whose moving-average coefficients are \(\Psi_h\). A structural representation has \(u_t=B\varepsilon_t\), \(BB'=\Sigma\), and orthogonal unit-variance innovations \(\varepsilon_t\). Require a positive output forecast-error variance at every evaluated horizon. Partition their labels into technology, demand, finance, and beliefs. Let \(\mathcal B\) be the matrices satisfying explicitly reported timing, sign, and external-instrument restrictions; belief innovations must exclude innovations in the specified observed fundamental-information set.
For output-selector \(e_Y\), define the horizon-\(H\) variance share of class \(k\) by
\[
\theta_k(B,H)=\frac{\sum_{h=0}^{H-1}\sum_{j\in k}(e_Y'\Psi_hB e_j)^2}
{\sum_{h=0}^{H-1}e_Y'\Psi_h\Sigma\Psi_h'e_Y}.
\]
The target is the joint identified set of these shares over \(B\in\mathcal B\), with sampling uncertainty and sensitivity to omitted information. Determine which additional independently justified measurements shrink that set enough to distinguish competing explanations of the main business-cycle component. A decomposition obtained from one selected rotation is insufficient: observationally equivalent rotations and alternative reasonable shock classifications must be reported. The exercise concerns a stated sample and model class, not universal fractions valid across all economies.

Take \(z_t\in\mathbb R^m\), \(\Sigma\succ0\), \(\Psi_0=I_m\), and a partition \((J_k)_{k=1}^4\) of \(\{1,\ldots,m\}\). The proposed shock classes are mutually exclusive labels, not automatically distinct economic mechanisms. Define \(\mathcal B(P)\) using the population observable law \(P\) and the specified restrictions; the target is \(\mathcal I_H(P)=\{(\theta_k(B,H))_{k=1}^4:B\in\mathcal B(P)\}\). Empty \(\mathcal B(P)\) means the restrictions are falsified. News is an innovation that revises a conditional forecast of a specified future fundamental; a belief shock is a separately defined subjective-forecast revision satisfying the declared exclusion. Orthogonality alone does not prove that a shock is nonfundamental.
\subsection*{Variants and assumption changes}
The following are \emph{editorial variants}, not additional published open conjectures.
\begin{enumerate}
\item \emph{Frequency-domain attribution.} For a fixed band \(D\subset[-\pi,\pi]\), replace forecast-horizon sums by integrated squared transfer functions over \(D\). Require an integrable spectral density and positive band variance; identify band shares separately from horizon shares.
\item \emph{External-instrument restriction.} Add observed \(Z_t\) with \(E[Z_t\varepsilon_{jt}]=0\) for every nontarget shock and nonzero covariance with the target. Compare the resulting identified set with sign-only restrictions, allowing bounded exclusion violations.
\end{enumerate}
\subsection*{References and source audit}
\begin{itemize}
\item George-Marios Angeletos; Fabrice Collard; Harris Dellas. \emph{Business-Cycle Anatomy}. 2020. Locator: Section III.C, final paragraph, printed p.3046; metadata on p.3030. Primary text checked. \url{https://economics.mit.edu/sites/default/files/publications/MBC_AER_published.pdf}.
\end{itemize}
The source leaves belief-versus-short-run-news interpretation unresolved; the four-way attribution is editorial.
\begin{enumerate}\raggedright
\item\hypertarget{definitions:source:30007512:1}{} George-Marios Angeletos; Fabrice Collard; Harris Dellas. 2020. \emph{Business-Cycle Anatomy}. Section III.C, final paragraph, printed p.3046; metadata on p.3030.
\url{https://economics.mit.edu/sites/default/files/publications/MBC_AER_published.pdf}.
DOI: \url{https://doi.org/10.1257/aer.20181174}.
\end{enumerate}

\subsection*{Common conventions: Source mathematical conventions}

These conventions apply to each entry unless explicitly replaced.

\textbf{Models and identification.} A primitive model \(m\) specifies all probability laws, feasible choices, information, equilibrium and measurement rules used by its target. The admissible class \(\mathcal M\) is fixed independently of outcomes. If \(P_O(m)\) is its observable probability law and \(\tau(m)\) the target, its identified set at observable law \(P\) is
\[
 \mathcal I_\tau(P)=\{\tau(m):m\in\mathcal M,\ P_O(m)=P\}.
\]
Point identification means that this set is a singleton. An empty set signals incompatibility of data and assumptions. Coordinate bounds need not describe the joint set. Bounds are sharp only if every retained target is attainable or arbitrarily approachable by compatible models. Finite bounds require explicit restrictions on unobserved quantities; unrestricted confounding, hidden wealth, extrapolation or arbitrary equilibrium selection can yield uninformative sets.

\textbf{Causal interventions.} A potential outcome \(Y_i(p)\) is the outcome under a complete policy regime \(p\), including timing, financing, information and market spillovers. The target population and units are fixed before treatment. With interference, use the complete assignment vector or a justified exposure mapping; individual no-interference notation alone cannot identify an economy-wide intervention. A difference of mean potential outcomes does not require the cross-world joint distribution. Conditioning on post-treatment migration, survival, enrollment or employment changes the estimand and needs separate justification.

\textbf{Statistical statements.} An estimator, confidence set, forecast or test specifies a sampling law, model class and loss/coverage criterion. Uniform \((1-\alpha)\)-coverage means \(\inf_{m\in\mathcal M}\Pr_m\{\tau(m)\in C_n\}\geq1-\alpha\), or an explicitly asymptotic version. The whole parameter space trivially has coverage, so informativeness requires a stated width, risk or power objective. A forecasting improvement is relative to a named benchmark, information set and evaluation population.

\textbf{Equilibrium and optimization.} An equilibrium comprises feasible plans, optimality, beliefs where needed, budget constraints and all market/resource clearing. The primitives or a stated benchmark define each correspondence \(\mathcal E(p;m)\). Nonemptiness is assumed or proved, never inferred just from compact policy sets. A selected-equilibrium target uses a declared selection rule; a robust target quantifies over all permitted equilibria. Use suprema or infima unless attainment is established. An \(\varepsilon\)-optimal policy is feasible and within \(\varepsilon\) of the finite optimum value. Report an empty feasible set separately.

\textbf{Welfare and accounting.} Normative utility, interpersonal normalization, social weights, numeraire, discounting and the use of public receipts are primitives. Behavioral utility need not coincide with the welfare criterion. Household welfare includes ownership income and rents once; adding profits again to compensating variation generally double-counts them. A monetary equivalent is defined by an explicit transfer and price convention, with monotonicity and range conditions. Average effects, mechanism contributions, transfers and welfare losses are different targets. Infinite sums require convergence and dynamic feasible plans require terminal or no-Ponzi conditions.

\textbf{Source versions.} A working-paper date, revision date and journal publication year are distinguished. A DOI for a journal version is not automatically the DOI of an earlier manuscript. Locators refer to the stated version and printed pages where specified. A metadata-only check does not verify a claim about a full-text conclusion. Every entry's source subsection narrows the provenance claim accordingly.
% END ECONOMICS STATEMENT
\end{document}
